\section{The object language and its three number-theoretic sentences}\label{sec:foundation}
The ordinary first-order presentation of intuitionistic primitive recursive arithmetic (iPRA) has one numerical sort, a symbol for each particular primitive recursive program, equality, intuitionistic connectives, and number quantifiers. Induction is restricted to quantifier-free formulas, which may include defined primitive recursive symbols. The repository presents the syntax as \code{PR k}, \code{PTm n} and \code{Fm PTm n}. A \code{PR k} is built from zero, successor, projections, composition and primitive recursion. Free variables of a term or formula lie among $n$ de Bruijn indices. The type \code{Prf PRAAx [] A} expresses that the closed sentence $A$ is derivable without hypotheses by the repository's rules.

Real numbers, functions quantified as objects, and sets of zeros do not occur in this language. We name the three hypotheses by arithmetic sentences, prove transformations inside iPRA where available, and identify their meaning in the standard model. A classical truth equivalence alone gives no term of the type \code{Prf PRAAx [] (Fm.iff A B)}. In particular, a classical argument producing a positive margin or a precision by existence has an additional obligation if one asks for a primitive recursive witness transformation.

\paragraph{Family parameters.}
Let $d_0,H,u,v$ be naturals. Write $d=d_0+1$, $b=H+2$ and $[w]_i=\lfloor w/b^i\rfloor\bmod b$. They name
\[
 F(X)=X^d+\sum_{i<d}([u]_i-[v]_i)X^i.
\]
Every monic nonconstant polynomial occurs by choosing a base above all coefficient magnitudes. A change of base can change the allowance below. A fixed code is retained when a bound depends on that allowance.

At each prime power put
\[
 h_F(p,k)=\deg\gcd(\bar F,X^{p^k}-X)\quad \text{over }\mathbb F_p,\qquad
 B_F(y)=\left\lfloor\log_2 \prod_{p^k\le y}p^{h_F(p,k)}\right\rfloor .
\]
Here $k\ge1$. The product of no factors is one. The gcd counts each distinct irreducible factor of $\bar F$ of degree dividing $k$, weighted by its degree. This follows by factoring $X^{p^k}-X$ over $\mathbb F_p$ and observing that its derivative is $-1$. The PR code checks divisibility and irreducibility by bounded searches through coefficient words, including when $\bar F$ has repeated roots. The position of the highest set bit, numbered from zero, implements the binary logarithm rounded down. Put
\[
 L(y)=1+\lfloor\log_2 y\rfloor\quad(y\ge1),\quad C=1000(d+1)^4(H+2)^2,\quad
 D_F(n)=B_F(2n)-2B_F(n).
\]
The original matrix is the excess
\[
 E(d_0,H,u,v,t)=\max\{0,D_F(n)^2-nC^2L(2n)^4\},\qquad n=100+t.
\]
RH is $\forall t\ E(0,0,0,1,t)=0$. GRH also quantifies $d_0$, with $(H,u,v)=(0,0,1)$. ERH quantifies all five arguments. With $F=X-1$, $B_F(y)$ is the index of the highest set bit of $\mathrm{lcm}(1,\ldots,\lfloor y\rfloor)$. These $\Pi^0_1$ formulas admit decidable finite failures. iPRA proves $\mathrm{ERH}\Rightarrow\mathrm{GRH}\Rightarrow\mathrm{RH}$ by universal specialization. For each it also proves equivalence to the negation of an existential failing tuple. In contrast, passing from the negation of a universal test to the existence of such a tuple would require an additional principle in general.

\paragraph{Why these integers recover the required functions.}
Factor $F=\prod_i f_i^{m_i}$ over $\Q$ with $f_i$ distinct and set $K_i=\Q[X]/f_i$. Let $G=\prod_i f_i$, $\Delta=|\mathrm{disc}(G)|$ and
$\Psi_F(y)=\sum_{p^k\le y}h_F(p,k)\log p$. Write $r$ for the number of factors. If $p\nmid\Delta$, reduction of the order generated by a root of $f_i$ gives the same finite algebra as reduction of its ring of integers. Distinct factors stay coprime and unramified. The residue degrees above $p$ therefore give exactly $h_F(p,k)$ when summed over degrees dividing $k$. At the exceptional primes either count is between $0$ and $d$. Consequently
\[
 \left|\Psi_F(y)-\sum_i\psi_{K_i}(y)\right|
 \le d\,\omega(\Delta)\log y,\qquad
 \sum_i\log |\Delta_{K_i}|\le\log\Delta\le 2d^2(H+2).
\]
The last inequalities follow from the discriminant of a power basis, the product formula for discriminants and resultants, and the root bound $|\alpha|\le H+2$. They hold in degree one with $\Delta=1$.

Under the zero hypothesis for the $K_i$, the explicit prime ideal theorem of Greni\'e and Molteni (Corollary 1.2) bounds $|\psi_K(y)-y|$ for $y\ge100$ by
\[
 \sqrt y\left[\left(\frac{\log y}{2\pi}+2\right)\log|\Delta_K|
          +\left(\frac{\log^2 y}{8\pi}+2\right)[K:\Q]\right].
\]
Adding over $i$ and using the exceptional-prime bound gives
$|\Psi_F(y)-ry|\le13d^3(H+2)\sqrt y\,L(y)^2$ at integer $y\ge100$.
Indeed each parenthesis is at most $3L(y)^2$ and $\omega(\Delta)\le2\log\Delta$.
The linear terms cancel at $2n$ and $n$. Flooring after division by $\log 2$ costs at most three bits, which is absorbed in $C\sqrt n L(2n)^2$. Thus every arithmetic test holds under the corresponding hypothesis. The converse follows already from the upper half alone by Theorem~\ref{thm:oneside}.

\paragraph{From the fields to the named hypotheses.}
Every number field has an integral primitive element, so its minimal polynomial is an allowed input. For $F=X^d-1$ the distinct factors are $\Phi_e$, $e\mid d$. The field $\Q(\zeta_e)$ has zeta function equal to the product of primitive Dirichlet $L$-functions of conductor dividing $e$, including conductor one. In the open strip all factors are holomorphic, so their zeros persist in the product with multiplicity. Taking $d$ equal to the conductor reaches each primitive character. The functional equations reflect nontrivial zeros across the critical line, for conjugate characters when necessary. Boundary zeros of imprimitive Euler factors on $\Re s=0$ have no role. For the terminology see the articles on the \href{https://en.wikipedia.org/wiki/Riemann_hypothesis}{Riemann hypothesis} and the \href{https://en.wikipedia.org/wiki/Generalized_Riemann_hypothesis}{generalized Riemann hypothesis}. For the analytic inputs see the \href{https://arxiv.org/abs/1312.4463}{explicit prime ideal theorem} and \href{https://math.mit.edu/classes/18.785/2018fa/LectureNotes19.pdf}{Sutherland, Theorem 19.15}.
